\section{Final-state-resolved response}\label{sec:8.7}

When a dark-matter particle deposits \SI{10}{eV} in silicon, it leaves an electron and a hole with about \SI{9}{eV} of kinetic energy between them. The split of the \SI{9}{eV} between the two carriers decides how many further pairs they make on their way down. The rate of impact ionization rises steeply with the energy of a carrier, so a lopsided split makes more pairs than an even one, and an even split is the one that minimizes impact ionization~\cite[Sect.~II\,A]{ramanathan2020ionization}. The step model of \cref{eq:1.5-step} assigns three pairs to every \SI{10}{eV} deposit, and the rate \cref{eq:8.2-spectrum} cannot do better, because the structure factor sums over all final states with the same $\omega$ and keeps no record of which one was made. The detector, however, counts the products of a particular final state. What must the response keep, so that the cascade of \crefrange{chap:18}{chap:20} receives the state the dark matter actually produced, and nothing is counted twice or lost on the way?

The answer is a response matrix resolved by the class of the final state, $S_{ab}(\vect{q},\omega;\Lambda)$, together with a rule for handing it over.

\subsection{Final-state classes}\label{sec:8.7-classes}

The cascade needs more than the sector of \cref{sec:8.6} in which a final state lies. It needs the state itself, at the moment the dark matter has finished acting and before the target starts to relax. Two time scales bracket that moment. Within a time of order $1/\delta_{\mathrm{F}}$, about a femtosecond for a pair in silicon, the electrons and phonons around the final state dress it; here $\delta_{\mathrm{F}}$ is the distance from the state to the nearest structure of the continuum it can decay into. The final state decays, by emitting its first phonon or by impact ionization, within $1/\varGamma_m$, with $\varGamma_m$ the width of \cref{sec:8.6}. Between the two times, the state is a well-defined \emph{prompt} final state. The interval between $1/\delta_{\mathrm{F}}$ and $1/\varGamma_m$ is the time window in which \cref{sec:8.6} found the split between $\mathcal{P}_0$ and $\mathcal{P}_+$ meaningful, and the Dyson convention \cref{eq:8.6-dyson} is the rule that labels a state inside it.

We group the prompt final states into \emph{final-state classes} $\Lambda$, each described by the quantum numbers the cascade needs. The commonest are these.
An electron--hole pair is labeled by the conduction band $c$ and crystal momentum $\vect{k}+\vect{q}$ of the electron and the valence band $v$ and crystal momentum $\vect{k}$ of the hole. These four labels fix the valleys the carriers occupy and their kinetic energies $E_{\mathrm{e}}$ and $E_{\mathrm{h}}$, measured from the band edges as in \cref{sec:6.5}. A spin label is added when the coupling acts on the spin. A pair with phonons carries in addition the branch, momentum and energy of each phonon. A core hole is labeled by its shell, since it relaxes by emitting an Auger electron (\cref{sec:10.6}). In the electron-following sector the class is the phonon content, or, for a freely recoiling nucleus, its energy $E_R$ and direction. Each class lies in one sector, so the classes refine the partition \cref{eq:8.6-partition} rather than cut across it.

The plasmon of silicon is not a class of its own. At the zero of $\varepsilon_1$ the measured $\varepsilon_2$ is 0.224 (\cref{we:4-siplasmon}), interband pairs exist at the plasmon energy, and the collective oscillation decays into them. In the random-phase approximation, the screened pair weights already contain the plasmon peak of the loss function, distributed over the pairs into which the plasmon decays (\cref{sec:8.7-pairs}). A plasmon needs a label only where it cannot decay into pairs, below the cutoff $q_{\mathrm{c}}$ of the electron gas of \cref{sec:4.5}, or where it can leave the target or convert into another excitation, as at the surfaces of films (\cref{sec:16.2}). It then has its own decay, which the cascade must treat.

For each class we write $P_\Lambda$ for the projector on its states. The classes are orthogonal and exhaust the final states, $P_\Lambda P_{\Lambda'} = \delta_{\Lambda\Lambda'}P_\Lambda$ and $\sum_\Lambda P_\Lambda = 1$. The response into one class follows from inserting $P_\Lambda$ around the energy delta function of \cref{eq:4.2-Sab}, as \cref{eq:8.6-sector} did for the sectors.
\begin{keyresult}[Final-state-resolved response]
\begin{equation}
  \begin{gathered}
  S_{ab}(\vect{q},\omega;\Lambda) = \frac{2\pi}{\Vol}\sum_ip_i\,\bigl\langle i\big|O_a(\vect{q})^\dagger\,P_\Lambda\,\delta(\omega + E_i - H)\,P_\Lambda\,O_b(\vect{q})\big|i\bigr\rangle , \\
  \sum_\Lambda S_{ab}(\vect{q},\omega;\Lambda) = S_{ab}(\vect{q},\omega) .
  \end{gathered}
  \label{eq:8.7-resolved}
\end{equation}
The sum rule holds exactly when each $P_\Lambda$ commutes with $H$, and each $S_{ab}(\vect{q},\omega;\Lambda)$ is then a positive semidefinite matrix.
\end{keyresult}
\noindent Both statements follow as for the sectors. When $P_\Lambda$ commutes with $H$, the delta function can be written in eigenstates of $H$ that each lie in one class, the cross terms between classes vanish, and each $S_{ab}(\vect{q},\omega;\Lambda)$ is a sum of outer products with positive weights. Since the prompt classes are defined on the window between $1/\delta_{\mathrm{F}}$ and $1/\varGamma_m$, they commute with $H$ only to the accuracy of that window. Inside the off-shell window of \cref{sec:8.6}, where the labels are unique to relative order $\theta_{\mathrm{ep}}$, each resolved response is positive. Near a zero-phonon resonance, the classes of $\mathcal{P}_+$ are not separately positive, and only their sum with the zero-phonon pair is, as \cref{eq:8.6-dyson} showed for the sectors. The delta part of the Dyson convention is a zero-phonon pair, weighted by its residue $z_{\mathrm{qp}}(m)$, and it enters $S_{ab}(\vect{q},\omega;\Lambda)$ as such. When this pair later emits a phonon, the emission is the first step of the cascade, not a second final state.

The rate into one class follows from the steps that gave \cref{eq:8.2-rate}, with $P_\Lambda$ inserted before and after the delta function. For one dark-matter particle of velocity $\vect{v}$,
\begin{equation}
  \varGamma(\vect{v};\Lambda) = \int\frac{\dd^3q}{(2\pi)^3}\sum_{a,b}\mathfrak{D}_{ab}(\vect{q},\vect{v})\,S_{ab}(\vect{q},\omega_{\vect{q}};\Lambda) ,
  \label{eq:8.7-classrate}
\end{equation}
and $\sum_\Lambda\varGamma(\vect{v};\Lambda) = \varGamma(\vect{v})$. For absorption, $\mathfrak{D}_{ab}$ is replaced by $g^{\phi\,*}_ag^\phi_b/2\omega_\phi$ at the single point $(\vect{p}_\phi,\omega_\phi)$, as in \cref{eq:8.2-abs}. In a crystal the resolved response depends on the direction of $\vect{q}$, as the total does, and so does the mix of classes a given particle produces.

\subsection{Pairs in an insulator}\label{sec:8.7-pairs}

For independent electrons the classes of the zero-phonon sector are single pairs. The ground state of a clamped insulator at $T = 0$ is the Slater determinant with every valence state filled. The density $n_e(\vect{q})$ of \cref{eq:3.2-ne} moves one electron from $v\vect{k}$ to $c,\vect{k}+\vect{q}$, with the amplitude $\rho_{cv}(\vect{k},\vect{q})$ of \cref{eq:3.2-rho} and the energy $\varepsilon_{c,\vect{k}+\vect{q}} - \varepsilon_{v\vect{k}}$. Each final state of \cref{eq:4.2-Sab} is therefore one pair. Summed over the two spins of the pair,
\begin{equation}
  S^0(\vect{q},\omega;cv\vect{k}) = \frac{4\pi}{\Vol}\,\bigl|\rho_{cv}(\vect{k},\vect{q})\bigr|^2\,\delta\bigl(\omega - \varepsilon_{c,\vect{k}+\vect{q}} + \varepsilon_{v\vect{k}}\bigr) .
  \label{eq:8.7-IP}
\end{equation}
\Cref{eq:8.7-IP} is the term of the independent-particle structure factor that belongs to one pair; summed over $c$, $v$ and $\vect{k}$ it is the crystal form of \cref{eq:2.5-golden}. Here $\vect{q}$ is not reduced to the first zone, and the folding vector $\vect{G}_{\mathrm{U}}$ of the final state stays in the vertex.

In the random-phase approximation, or in time-dependent density-functional theory, the external potential at $\vect{q}$ induces a density, and the density adds its own Hartree and exchange--correlation potentials at every $\vect{q} + \vect{G}$. The electrons, independent in the Kohn--Sham sense, respond to the total potential. The total potential is the external potential times the column
\begin{equation*}
  \bigl(\tilde\varepsilon^{-1}\bigr)_{\vect{G}0}(\vect{q},\omega) = \delta_{\vect{G}0} + \sum_{\vect{G}'}\bigl(f_{\mathrm{Hxc}}\bigr)_{\vect{G}\vect{G}'}(\vect{q},\omega)\,R_{\vect{G}'0}(\vect{q},\omega) ,
\end{equation*}
the inverse dielectric matrix of a test electron, which \cref{eq:6.2-screen} met for the static potential of a phonon. Without the exchange--correlation kernel it is the inverse dielectric matrix \cref{eq:4.4-epsGG}. The index $0$ is the row $n_e(\vect{q})$ itself, and $R$ is the retarded response. It is natural to guess that the screened pair weight replaces $\rho_{cv}(\vect{k},\vect{q})$ in \cref{eq:8.7-IP} by the vertex of the total potential, $\sum_{\vect{G}}(\tilde\varepsilon^{-1})_{\vect{G}0}\,\rho_{cv}(\vect{k},\vect{q}+\vect{G})$. The guess is right, and the algebra that confirms it admits a second answer.

Let $f = f_{\mathrm{Hxc}}$ be Hermitian as a matrix in the reciprocal-lattice vectors, as it is in the random-phase approximation and for any real kernel without frequency dependence, such as the adiabatic local-density kernel of \cref{sec:4.6}. Write $R^\dagger$ for the Hermitian conjugate matrix, the advanced response. The Dyson equation \cref{eq:4.6-tddft} reads $R = R^0 + R^0fR$, and its conjugate $R^\dagger = R^{0\dagger} + R^\dagger fR^{0\dagger}$. The Dyson equation gives $R^0(1 + fR) = R$, its conjugate gives $(1 + R^\dagger f)R^{0\dagger} = R^\dagger$, and together
\begin{equation}
  (1 + R^\dagger f)\,\bigl(R^0 - R^{0\dagger}\bigr)\,(1 + fR) = (R + R^\dagger fR) - (R^\dagger + R^\dagger fR) = R - R^\dagger ,
  \label{eq:8.7-factor}
\end{equation}
for every target. The fluctuation--dissipation theorem \cref{eq:4.2-fdt} turns both differences into structure factors: at $T = 0$ and $\omega > 0$, $S_{\vect{G}\vect{G}'} = \ii(R - R^\dagger)_{\vect{G}\vect{G}'}$, and the same holds for $R^0$. In the bubble \cref{eq:4.6-R0} only the resonant term has an imaginary part at $\omega > 0$, and $\ii[1/(x + \ii0^+) - 1/(x - \ii0^+)] = 2\pi\delta(x)$, so
\begin{equation*}
  \ii\bigl(R^0 - R^{0\dagger}\bigr)_{\vect{G}\vect{G}'} = \frac{4\pi}{\Vol}\sum_{cv\vect{k}}\rho_{cv}(\vect{k},\vect{q}+\vect{G})^*\,\rho_{cv}(\vect{k},\vect{q}+\vect{G}')\,\delta\bigl(\omega - \varepsilon_{c,\vect{k}+\vect{q}} + \varepsilon_{v\vect{k}}\bigr) .
\end{equation*}
Take the element $00$ of \cref{eq:8.7-factor}. Since $f$ is Hermitian, $(1 + R^\dagger f)_{0\vect{G}} = [(1 + fR)_{\vect{G}0}]^* = [(\tilde\varepsilon^{-1})_{\vect{G}0}]^*$, and the double sum over $\vect{G}$ and $\vect{G}'$ closes into a squared modulus for each pair.
\begin{keyresult}[Screened pair-resolved response]
\begin{equation}
  S(\vect{q},\omega;cv\vect{k}) = \frac{4\pi}{\Vol}\,\Bigl|\sum_{\vect{G}}\bigl(\tilde\varepsilon^{-1}\bigr)_{\vect{G}0}(\vect{q},\omega)\,\rho_{cv}(\vect{k},\vect{q}+\vect{G})\Bigr|^2\,\delta\bigl(\omega - \varepsilon_{c,\vect{k}+\vect{q}} + \varepsilon_{v\vect{k}}\bigr) ,
  \label{eq:8.7-screened}
\end{equation}
with the retarded $\tilde\varepsilon^{-1}$, in the random-phase approximation and in time-dependent density-functional theory with a Hermitian kernel. The sum over $c$, $v$ and $\vect{k}$ is the screened structure factor $S(\vect{q},\omega) = (q^2/2\pi\alpha)L(\vect{q},\omega)$ of \cref{eq:4.4-SL}, local fields included.
\end{keyresult}
\noindent Each term is positive, the terms add up to the full screened response, and the plasmon peak of $L$ is in the sum, spread over the pairs into which the plasmon decays. The pair weight is the bare vertex of the total potential, evaluated at the energy deposited. In the deviation notation of \cref{sec:8.5}, \cref{eq:8.7-IP} is at $\mathcal{X} = \mathrm{IP}$ and \cref{eq:8.7-screened} at $\mathcal{X} = \mathrm{RPA} + \mathrm{LF}$ or TDDFT; with the BSE the final states are eigenstates of the pair Hamiltonian, excitons among them, and these eigenstates become the classes (\cref{sec:10.5}).

The second answer comes from the same algebra in the other order. The Dyson equation can equally be written $R = R^0 + RfR^0$, with conjugate $R^\dagger = R^{0\dagger} + R^{0\dagger}fR^\dagger$, and then the identity $(1 + Rf)(R^0 - R^{0\dagger})(1 + fR^\dagger) = R - R^\dagger$ holds as well, and it distributes the same total over the pairs with the advanced weight $\delta_{\vect{G}0} + \sum_{\vect{G}'}f_{\vect{G}\vect{G}'}R^*_{0\vect{G}'}$ in place of $(\tilde\varepsilon^{-1})_{\vect{G}0}$. The two distributions have the same sum for every target, and pair by pair they need not agree. Time decides which of the two distributions labels the final state. The pair that the cascade receives is the state found at late times, an \emph{out-state} of scattering theory, and the retarded resolvent enters through the overlap of a free final state with the evolving one: in Lippmann--Schwinger form, $\langle\psi_f^{(-)}| = \langle\phi_f|(1 + VG^{(+)})$~\cite[Chap.~6]{sakurai2017modern}. In the time-dependent mean-field picture the choice of the retarded weight is plain causality. The external potential acts first and induces the total potential $(\tilde\varepsilon^{-1})_{\vect{G}0}\phi_{\mathrm{ext}}$; the total potential creates bare pairs, and the bare pairs are the asymptotic states of the mean-field evolution. The advanced form labels each pair by the state it is evolved from, an in-state, and the in-state is not the label the cascade needs.

The two labelings agree pair by pair whenever $R(\vect{q},\omega)$ is Hermitian, since then the two weights coincide. Inside the off-shell window $R$ is Hermitian up to terms of order $g^2$: the zero-phonon pairs are closed, and the anti-Hermitian part of the external response comes from the phonon-assisted absorption itself; the difference in a pair-resolved rate is then of order $g^4$. The two labelings also agree in every target that has an even number of electrons and is invariant under the product $\mathsf{P}\mathsf{T}$ of inversion and time reversal, for the reason that made the Fano coefficients of \cref{eq:8.6-fano} vanish. The $\mathsf{P}\mathsf{T}$ argument needs a kernel that is real symmetric, $f_{\vect{G}\vect{G}'} = f_{\vect{G}'\vect{G}} = f^*_{\vect{G}\vect{G}'}$, as the random-phase and adiabatic local-density kernels are; a kernel that is merely Hermitian does not suffice. Every state can then be chosen invariant under $\mathsf{P}\mathsf{T}$, and with the origin at an inversion center $n_e(\vect{q}+\vect{G})$ is invariant too, since inversion reverses $\vect{q}+\vect{G}$ and complex conjugation restores it. Every amplitude $\langle f|n_e(\vect{q}+\vect{G})|i\rangle$ is then real, $R_{\vect{G}\vect{G}'} = R_{\vect{G}'\vect{G}}$, and the advanced weight is the complex conjugate of the retarded one. The symmetry of $R$ depends on that choice of origin, since a shift multiplies $R_{\vect{G}\vect{G}'}$ by phases; the pair weights do not. A real vertex sum has the same modulus with the retarded and with the advanced weight, so each pair, or each set of degenerate pairs, gets the same weight. Silicon and germanium satisfy the hypothesis. GaAs has no center of inversion, and there the choice of the out-state matters.

The pair weights need nothing beyond the ingredients of the structure factor itself. The vertices $\rho_{cv}(\vect{k},\vect{q}+\vect{G})$ come from the plane-wave coefficients of the Bloch states (\cref{app:E.7}), and the column $(\tilde\varepsilon^{-1})_{\vect{G}0}$ from a dielectric-matrix calculation at the same $(\vect{q},\omega)$. A structure-factor calculation and a resolved calculation differ in what they keep. A code that returns $S(\vect{q},\omega)$ sums the pairs and discards them. The resolved response stores their weights, binned in the carrier energies and valleys, for each $\vect{q}$, $\omega$ and coupling row, and a response library must say which binning it keeps.

\begin{workedexample}{Pairs made at one momentum transfer in silicon}\label{we:8-pairs}
Which pairs does a dark-matter kick make, and how do they differ from the pairs that a photon of the same energy makes? We take the empirical pseudopotential of Cohen and Bergstresser for silicon~\cite{cohen1966band}, the model of \cref{sec:3.3,sec:4.6}: three form factors fitted to measured optical transitions, $a = \SI{5.43}{\angstrom}$, 113 plane waves and no spin--orbit coupling. Its indirect gap is \SI{0.82}{eV}, below the measured \SI{1.17}{eV}, with the conduction minimum at $0.853$ of the way to X. The momentum transfer is $q_0 = \alpha m_e = \SI{3.73}{keV}$ along $[100]$, a typical value for electron scattering, which is $1.633\,(2\pi/a)$ and lies beyond the first zone; its Umklapp partner $\vect{q} + \vect{G}$ with $\vect{G} = (-2,0,0)\,2\pi/a$ has $|\vect{q}+\vect{G}| = \SI{0.84}{keV}$. The energy deposited is $\omega = \SI{10}{eV}$. The sums run over a shifted $24^3$ grid of $\vect{k}$ over the full zone, with 4 valence and 26 conduction bands. The local-field matrix has the 51 vectors with $|\vect{q}+\vect{G}|^2 \le 12\,(2\pi/a)^2$, and the $\ii0^+$ of the bubble is a Lorentzian of half-width \SI{0.10}{eV}. The energy delta function of each pair is a Gaussian of standard deviation \SI{0.10}{eV} instead. A Lorentzian has long tails, and in the optical limit below, whose weights carry $1/(\varepsilon_c - \varepsilon_v)^2$, the tails would admit strong transitions several eV below $\omega$ and break the energy balance of the pairs. The numbers quoted change by at most \SI{3}{\percent} between $20^3$ and $24^3$ points and between Gaussians of 0.10 and \SI{0.05}{eV}.

\emph{Screening.} The independent-particle structure factor is $S^{\mathrm{IP}}(\vect{q}_0,\omega) = \SI{1.7e8}{eV^2}$. The random-phase approximation with local fields gives $(\varepsilon^{-1})_{00} = 0.55 - 0.16\,\ii$, and summing \cref{eq:8.7-screened} over pairs gives the screened structure factor \SI{4.9e7}{eV^2}, a factor of 0.29 of the unscreened one. The head of the column does most of the screening, since $|(\varepsilon^{-1})_{00}|^2 = 0.33$; the other 50 elements are at most 0.05 each, the Umklapp partner among the small ones, and together they lower the rate by a further \SI{12}{\percent}. For comparison, the free gas at the valence density gives $m_e^2\omega/\pi q_0 = \SI{2.2e8}{eV^2}$ at the same point, the unpolarized form of the one-spin structure factor of \cref{we:8-cross}.

\emph{Which pairs.} Every pair at this $\omega$ shares $\omega - \Eg = \SI{9.2}{eV}$ between the kinetic energies $E_{\mathrm{e}}$ and $E_{\mathrm{h}}$ of its carriers. With the screened weights, the mean electron energy is \SI{4.8}{eV} and the mean hole energy \SI{4.4}{eV}; the electron takes more than half in \SI{64}{\percent} of the weight, and one carrier takes more than three quarters in \SI{45}{\percent}. The lowest valence band supplies \SI{4}{\percent} of the holes and the other three about a third each. The independent-particle weights give the same picture, each of these fractions within three percentage points. Screening lowers the number of pairs made by a factor of three and leaves their distribution nearly unchanged.

\emph{Other seeds at the same energy.} A photon of \SI{10}{eV}, or a dark photon, makes vertical pairs with the weights $|\hat{\vect{e}}\cdot\vect{v}_{cv}|^2/(\varepsilon_c - \varepsilon_v)^2$ of the optical limit \cref{eq:3.2-optical}. Its electrons are more energetic, \SI{5.7}{eV} on average, and its holes less, \SI{3.5}{eV}, about \SI{0.9}{eV} below the holes of the dark-matter kick; the electron takes more than half in \SI{85}{\percent} of the weight. The transverse current $\hat{\vect{y}}\cdot\vect{j}_e(\vect{q}_0)$ at the same momentum transfer lies in between, with \SI{5.3}{eV} and \SI{3.8}{eV}, and puts half its weight on lopsided pairs with one carrier above three quarters. Every seed shares the same \SI{9.2}{eV}, and yet the three seeds hand the cascade three different starting distributions.

\emph{In-states and out-states.} In this model of silicon, with the origin at the bond center, an inversion center, the two labelings give every pair the same weight, and $R_{\vect{G}\vect{G}'} = R_{\vect{G}'\vect{G}}$, as the $\mathsf{P}\mathsf{T}$ argument requires for the real symmetric kernel of the random-phase approximation. The same model of GaAs has no center of inversion. At the same $q_0$ and \SI{10}{eV}, the pair-by-pair differences between the two labelings add up, in absolute value, to \SI{15}{\percent} of the screened total, although the mean carrier energies of the two labelings differ by \SI{0.02}{eV} at most. The choice of the out-state matters for momentum- and valley-resolved seeds, and much less for the energies alone.
\end{workedexample}

\subsection{Cascade probabilities of a final state}\label{sec:8.7-handoff}

Once the dark matter has acted, the target evolves under its own Hamiltonian. A pair of class $\Lambda$ emits phonons, ionizes further pairs and slows down, and the event ends with some number $N_{\mathrm{eh}}$ of pairs, some energy in phonons and some energy stored in defects. We write $P_{\mathrm{casc}}(N_{\mathrm{eh}}|\Lambda)$ for the probability that the cascade started by a final state of class $\Lambda$ ends with $N_{\mathrm{eh}}$ pairs; \cref{sec:18.3} computes it from a branching process. We call \emph{seeding} the step that turns the prompt state into the initial condition of the cascade. The one assumption here is that the class carries everything the later evolution depends on. The prompt state is reached before the first decay, so the history of how it was made can enter only through coherences between production and decay. These are the Fano coefficients of \cref{eq:8.6-fano}, and they vanish in silicon and germanium.

Under that assumption, the probability of an event with $N_{\mathrm{eh}}$ pairs is the probability of each class times the probability that its cascade ends with $N_{\mathrm{eh}}$ pairs, summed over the classes. With the rate into a class, \cref{eq:8.7-classrate}, and the halo average of \cref{eq:8.2-spectrum}:
\begin{keyresult}[Hand-off to the cascade]
\begin{equation}
  \mathcal{R}_{N_{\mathrm{eh}}} = \frac{\rho_\chi}{\mchi\,\rho_T}\int\dd^3v\,f_\chi(\vect{v})\int\frac{\dd^3q}{(2\pi)^3}\sum_\Lambda\sum_{a,b}\mathfrak{D}_{ab}(\vect{q},\vect{v})\,S_{ab}(\vect{q},\omega_{\vect{q}};\Lambda)\,P_{\mathrm{casc}}(N_{\mathrm{eh}}|\Lambda) ,
  \label{eq:8.7-handoff}
\end{equation}
the rate per unit target mass of events with $N_{\mathrm{eh}}$ pairs. The coupling tensor is contracted inside the sum over classes, and the sum over classes is inside the integrals over $\vect{q}$ and $\vect{v}$. Equivalently, at fixed $(\vect{q},\omega,\vect{v})$ the pair number is distributed as
\begin{equation}
  P(N_{\mathrm{eh}}|\vect{q},\omega,\vect{v}) = \frac{\sum_\Lambda\sum_{a,b}\mathfrak{D}_{ab}(\vect{q},\vect{v})\,S_{ab}(\vect{q},\omega;\Lambda)\,P_{\mathrm{casc}}(N_{\mathrm{eh}}|\Lambda)}{\sum_{a,b}\mathfrak{D}_{ab}(\vect{q},\vect{v})\,S_{ab}(\vect{q},\omega)} .
  \label{eq:8.7-cond}
\end{equation}
\end{keyresult}
\noindent The same form holds for any other count of the cascade, the phonon energy or the energy stored in defects, with $P_{\mathrm{casc}}$ replaced by the distribution of that quantity, and for absorption, with $\mathfrak{D}_{ab}$ replaced as in \cref{eq:8.7-classrate}. \Cref{eq:8.7-handoff} contains the pair bins of \cref{eq:1.5-bins} as a special case. If $P_{\mathrm{casc}}(N_{\mathrm{eh}}|\Lambda)$ depended on the class only through its energy $\omega$, \cref{eq:8.7-cond} would reduce to $P_{\mathrm{casc}}(N_{\mathrm{eh}}|\omega)$, the sum over $\Lambda$ would close into $S_{ab}(\vect{q},\omega)$ by \cref{eq:8.7-resolved}, and the integrals would collapse to $\int\dd\omega\,(\dd\mathcal{R}/\dd\omega)\,P_{\mathrm{casc}}(N_{\mathrm{eh}}|\omega)$. The step model is the further choice of a distribution that puts all its weight on $1 + \lfloor(\omega - \Eg)/\varepsilon_{\mathrm{eh}}\rfloor$.

That reduction is the common shortcut, and \cref{eq:8.7-handoff} says when it fails.

\emph{The momentum transfer stays inside.} At one $\omega$, different $\vect{q}$ make different pairs. \Cref{we:8-pairs} shows that at \SI{10}{eV} in silicon the hole carries about \SI{0.9}{eV} more, on average, for a dark-matter kick at $q_0$ than for an absorbed photon. This difference of \SI{0.9}{eV} is close to the gap, the energy scale on which impact ionization switches on. If the spectrum is integrated over $\vect{q}$ first and the cascade is seeded afterwards with a distribution of pairs chosen at fixed $\omega$, the pairs of one momentum are attached to the rate of another. The resolved response depends on the direction of $\vect{q}$ as well, so in a crystal the mix of seeds follows the dark-matter wind through the day, and a directional analysis needs the seeds at every direction it uses.

\emph{The velocity and the model stay inside.} With one coupling row $a$, the weights of the classes in \cref{eq:8.7-cond} are $S_{aa}(\vect{q},\omega;\Lambda)/S_{aa}(\vect{q},\omega)$, independent of $\vect{v}$ and of the model, and a table of $P(N_{\mathrm{eh}}|\vect{q},\omega)$ per row serves every model that couples through that row. With several rows a table per row does not serve every model. The weights then depend on $\mathfrak{D}_{ab}(\vect{q},\vect{v})$, so on the coefficients of the model and on the velocity, and the off-diagonal elements $S_{ab}(\vect{q},\omega;\Lambda)$ populate the classes too. If single-row tables are mixed in proportion to the single-row rates, the mixture misses the interference of \cref{sec:8.2-cross} class by class, even where the interference cancels in the total.

\emph{The operator structure stays inside.} Different rows probe different parts of the electron's state and make different pairs at the same $(\vect{q},\omega)$. A spin row reaches spin-flipped pairs that the density cannot. A current row weights each pair by the velocity of its carriers, not by the overlap of their periodic parts. The longitudinal current is the exception: the continuity equation \cref{eq:4.3-continuity} gives $\langle f|\hat{\vect{q}}\cdot\vect{j}_e(\vect{q})|i\rangle = (\omega/q)\langle f|n_e(\vect{q})|i\rangle$ for every pair of eigenstates, so the longitudinal current makes exactly the pairs the density makes. The transverse current, the spin density and the spin current make their own pairs, and the velocity-dependent operators of \cref{tab:8-operators} reach the target through them.

The cascade depends on $\vect{q}$ mainly through its first generation. The electron and hole of the prompt pair carry the crystal momenta $\vect{k}+\vect{q}$ and $\vect{k}$. Where collisions randomize the momentum of a carrier faster than they change its energy, its later history depends on its energy alone, and the secondary carriers inherit nothing of $\vect{q}$; \cref{sec:18.1} states this reduction and the energies at which it fails. The difference between dark-matter and photon seeds then enters once, through the distribution of $(E_{\mathrm{e}},E_{\mathrm{h}})$ and of the valleys at the start. This one difference is enough to change the pair count at the few-pair level, where a single impact ionization more or less is a whole bin. The ionization model of Ramanathan and Kurinsky describes the split of the energy between the carriers by one distribution of the electron's energy at each total, fitted to the ionization yields of photon absorption~\cite[Eqs.~(5)--(7)]{ramanathan2020ionization}. This distribution is the $\vect{q} \to 0$ marginal of the seeds of \cref{eq:8.7-cond}, and using it for dark-matter scattering replaces the seeds of one momentum by the seeds of another; \cref{sec:18.3} sizes the error. For absorption the momentum is $p_\phi = m_\phi v \approx 10^{-3}m_\phi$, the transitions are vertical, and a dark photon is absorbed through the same current as a photon. Its seeds are the photon's, and a photon calibration applies to it exactly.

We have derived the identification of the prompt pair with the out-state of \cref{eq:8.7-screened} only within time-dependent mean-field linear response, where the bare pair is the asymptotic state of the evolution. The exact final states of an interacting insulator are not labeled by pairs, and it is not established whether the retarded weights give the distribution of the carriers found at late times beyond mean field. The assumption behind \cref{eq:8.7-handoff} is that the class fixes the statistics of the cascade. This assumption holds when the only coherence between production and decay is the one the Fano coefficients carry; the widths of the prompt states would set a bound on what else the history of an event can contribute, but this bound is not known.
